% !PRINTED 782
\weberpage{782}
Andererseits ist aber nach 2.:

\weq{19}{(s-1) Q_1 Q_2 \ldots Q_\mu}

für $s = 1$ endlich, und wenn wir also (19) durch (18) dividiren, so folgt:
\[
(s-1)^{1-\frac{\mu}{\nu}} \text{ für } s = 1 \text{ endlich,}
\]
d.\,h.

\weq{20}{\nu \geqq \mu.}

Dies mit (14) zusammen ergiebt aber $\nu = \mu$, wodurch der folgende Satz bewiesen ist:

3. \gesperrt{Der Grad des Kreistheilungskörpers $\Omega_m$ ist gleich $\varphi(m)$, also die Kreistheilungsgleichung $\varphi(m)^{\mathrm{ten}}$ Grades, deren Wurzeln die primitiven $m^{\mathrm{ten}}$ Einheitswurzeln sind, irreducibel.}

Diese Betrachtung leistet uns aber noch einen anderen wichtigen Dienst, indem sie uns den Beweis des berühmten Satzes liefert, \gesperrt{dass in jeder arithmetischen Progression, deren Differenz und Anfangsglied natürliche Zahlen ohne gemeinsamen Theiler sind, unendlich viele Primzahlen enthalten sind}\textsuperscript{1)}.

Es folgt nämlich jetzt aus (18), dass das Product
\[
(s-1) Q_1 Q_2 \ldots Q_\mu
\]
für $s = 1$ einen endlichen und von Null verschiedenen Werth hat, und da nach 2. keiner der Factoren unendlich ist, so folgt:

4. \gesperrt{Die Summen}
\[
(s-1) Q_1(s), Q_2(s), \ldots, Q_\mu(s)
\]
\gesperrt{haben für $s = 1$ endliche und von Null verschiedene Werthe.}

Wenden wir auf alle Factoren von (11) die Formel an:
\[
\log \left[1 - \chi(p) p^{-s}\right] = \frac{\chi(p)}{p^s} + \frac{1}{2} \frac{\chi(p^2)}{p^{2s}} + \frac{1}{3} \frac{\chi(p^3)}{p^{3s}} + \cdots % [sic] im Druck fehlt das mathematisch erforderliche Minuszeichen vor \log
\]

\noindent\rule{4cm}{0.4pt}

\textsuperscript{1)} Der erste vollständige und allgemeine Beweis dieses Satzes ist von \gesperrt{Dirichlet} gegeben (Abhandlungen der Berl. Akademie vom Jahre 1837. Gesammelte Werke Nr.\ XXI). Auf die wesentliche Vereinfachung dieses Beweises durch Benutzung der \gesperrt{Kummer}'schen Formeln für die Classenzahl in den Kreistheilungskörpern hat \gesperrt{Dedekind} aufmerksam gemacht (Vorlesungen über Zahlentheorie, 3. Auflage, S. 596).
%CONT>>
